\section{Separation of packets}\label{sec:separation}


We expand an analytic function successively along a tree: each
coefficient at one level is expanded again at the next.  No convergence
of a series in all the quantities
$h_i=\exp(-Z_i)$ is asserted or used.  The distinction is particularly
important when a coefficient has zero formal data: decay to every fixed
order at its own level must first be strengthened before that coefficient
can be passed at the preceding level.

Uniqueness of an analytic function from its complete asymptotic data
belongs to the quasianalytic approach developed in the
limit-cycle literature \citep{Ilyashenko1991,Ecalle1992}.
Transserial Ilyashenko algebras provide a related injectivity theorem
for their own expansion classes \citep{GalalKaiserSpeissegger2020}.
Here the hypotheses are the precise band, transition, and source
estimates below, and the proof establishes the needed implication
directly for those data.

\subsection{Data and quantifiers}

Fix a saturated flag and the paths of Lemma~\ref{lem:flags-path}.
We use $s=u+iv$, $f_i=\log Z_i$, $\Lambda_i=f_i'$ on the real
path, and
\[
 V_i(u)=\frac{e^{f_i(u)}}{\Lambda_i(u)}.
\]
All logarithms have their lifted determinations.  The symbol $b_j$
below denotes the entire vector of bounded backgrounds in band $j$;
its components enter the triangular flag equations.  Accordingly,
$f_{k,j}$ and $Z_{k,j}$ denote the $k$th logarithm and scale evaluated with
that vector.  The subscripts on $b_j$ do not denote the index of an
individual triangular equation.

A \emph{good contour of level $i$} means either of the contours
of Lemma~\ref{lem:flags-contours} with
\begin{equation}\label{eq:separation-good}
 \cos(\Im f_i)=\frac{\epsilon_i}{\Lambda_i},\qquad
 V_i=\frac{e^{f_i}}{\Lambda_i},
\end{equation}
where $\epsilon_i>0$ is fixed and can be chosen arbitrarily small.
Contours always have small collars.  A \emph{side of level $i$}
is the boundary of the positive-real-part region of a holomorphic
proxy for
\begin{equation}\label{eq:separation-barrier}
 L_i=Z_i-A\sum_{k>i}Z_k-\frac{Z_i}{f_i^K}.
\end{equation}
Here $A$ is sufficiently large for the finite calculation under
consideration, and $K$ lies in a fixed interval with room to decrease
it finitely, for example $2\leq K\leq4$.  Proxy construction is proved
below.  Before constructing the proxy, the domain has extra angular
room.  A later increase of $A$ and decrease of $K$ puts the final side
strictly within that room.

The exceptional constants $D_i$ are those of
Lemma~\ref{lem:flags-exceptional}: $D_i=\log(2+x_0)$ when $f_i$
contains an uncapped slow summand with coefficient $x_0$, and $D_i=1$
otherwise; put $D_0=1$.  These constants remain fixed on each chosen
path.  They need not remain bounded as its anchor tends to infinity.

\begin{definition}[Packet data]\label{def:packet}
Let $t_c$ be a real ordinary germ point.  A packet consists of the
following data and estimates, separately for signs $+$ and $-$.

\emph{Tree and actual function.}
For every prefix $\alpha=(a_1,\ldots,a_{j-1})$ there is a raw
function $S^{j,\pm}_{\alpha}$ in band $j$, and a set
$E^{j,\pm}_{\alpha}\subset\RR$ of child exponents.  Each such set
is finite below every finite ceiling.  At level $m+1$ the functions
$c^{\pm}_\alpha=S^{m+1,\pm}_\alpha$ are analytic ordinary germs
at $t_c$.  The two functions at level $1$ are restrictions of one
actual holomorphic function across the real ray.  Absent exponents
can be inserted with identically zero packets.  Thus the supports
are \emph{iterated left-finite}, including at prefixes whose earlier
coordinates need not be positive.

\emph{Independent tests.}
The estimates are required along the independent free paths of
Section~\ref{sec:flags}, with independent ordinary tests $t(s)$.
These tests are holomorphic, real on the ray, bounded in their
ordinary neighborhoods, and have positive-order derivatives tending
to zero through every fixed requested order.  They can approach any
specified ordinary value on a fixed window about a varying anchor,
and return to $t_c$, with derivatives tending to zero, at infinity
on every individual path.  Backgrounds have the controlled real
Taylor comparisons of Lemma~\ref{lem:flags-background}.  Whenever
a determination's own thin side must be located, the comparisons
are made to sufficiently high fixed Taylor order.

\emph{Finite requests and uniformity.}
The following requirements apply to every finite choice of labels,
normalizations, budgets, and backsteps.  Ordinary neighborhoods may
be shrunk and domains narrowed a finite number of times.  Constants
and sufficiently far starts may depend on this finite request.
For a fixed such request they are uniform as anchors recede, on
the full forward domains of controlled test families with common
ordinary bounds and slack, common bounds for the required finite
jets and Taylor errors, and common moduli in the required small-jet
and flag comparisons.  In particular the finite source counts,
prefactors, positive precision constants, and eventual absorption
of the stated $o(U)$ losses are uniform for that family.  The
allowed exceptional constants $D_i$ are retained explicitly in
the growth bounds.  This is uniformity for one fixed finite
request; it is not uniformity over all truncation orders.
Previously chosen inner functions agree on overlaps with the ones
used in later requests.  All further fixed-budget estimates hold
eventually at infinity on the \emph{same} path, even if their
ordinary neighborhoods are smaller and their starts are farther out.
The return time to a smaller ordinary neighborhood may depend on
the test.  There is no uniform start over all budgets or all such
return times.

\begin{enumerate}
\item \emph{Bands and backgrounds.}
Band $j$ extends from any sufficiently vertical preceding
good contour, with a collar, to its own side
\eqref{eq:separation-barrier}.  For $j=1$ its inner boundary is the
real ray; its data are genuinely holomorphic across that ray and
have no antiholomorphic sources.  Each band has a controlled
background $b_j$.  At a selected level-$j$ good collar,
\begin{equation}\label{eq:separation-background-join}
 b_j-b_{j+1}=O(e^{-cV_j}).
\end{equation}
For leading estimates through the entire outer fringe we additionally
require, after narrowing,
\begin{equation}\label{eq:separation-background-fringe}
 b_j-b_{j+1}=O(e^{-a\Re Z_{j,j}}),\qquad a>0.
\end{equation}
Both estimates compare determinations at the same point.

\item \emph{Transitions.}
On the outer fringe of band $j$, from sufficiently vertical
good collars to its side, for every finite real $M$ there is a finite
cutoff $M'>M$ such that
\begin{equation}\label{eq:separation-transition}
 S^{j,\pm}_{\alpha}
  =\sum_{\substack{a\in E^{j,\pm}_{\alpha}\\a<M'}}
     e^{-aZ_{j,j}}S^{j+1,\pm}_{\alpha,a}
       +O(e^{-M\Re Z_{j,j}}).
\end{equation}
The fringe satisfies
$\cos(\Im f_{j,j})\leq\epsilon/\Lambda_j$ for sufficiently
small fixed $\epsilon$, and $A$ in
\eqref{eq:separation-barrier} can be enlarged for this budget.
New contour choices at this or later transitions do not alter the
already selected inner joins.

\item \emph{Growth and safe columns.}
Throughout band $j$,
\begin{equation}\label{eq:separation-raw-growth}
 \log^+|S^{j,\pm}_{\alpha}|
       \leq C\Re Z_{j,j}+CD_{j-1}.
\end{equation}
The $D_{j-1}$ term is absent near the band's own side and on safe
vertical columns.  For every finite set of raw labels and bands,
there are simultaneous safe columns in every fixed-length real
interval sufficiently far along the path.  If
\[
 D_{j-1}>\frac{Z_j(u_*)}{f_j(u_*)^{K+3}},
\]
the same improvement holds on the entering level-$(j-1)$ good
collars for
\[
 |f_{j-1}(u)-f_{j-1}(u_*)|
       \leq C_1\log(2+D_{j-1}),
\]
for every fixed $C_1$ needed in the calculation.  The constants
are uniform on these safe sets.  For a controlled family the
columns themselves may depend on its member; simultaneity
concerns the finite labels and bands within that member.  The
search-interval lengths and all bounds are uniform in the family.

\item \emph{Antiholomorphic sources.}
In bands $j>1$ the data may be $C^1$ interpolations of local
holomorphic determinations.  Their $\bar\partial_s$ derivatives
are bounded by finite sums of $e^{-cU}$, with sources supported
where those determinations or cutoffs change.  A permitted cutoff
cost in band $j$, with $r$ a dependent node, is
\begin{equation}\label{eq:separation-costs}
 U=e^{\Re f_{r,j}}\quad\hbox{or}\quad
 U=e^{\Re f_{r,j}}\cos(\Im f_{k,j}),
 \qquad r<j,\quad\sigma(r)\geq j,\quad k\geq j,
\end{equation}
or any cost bounded below by a fixed positive multiple of one of
these.  Smoothing a level-$i$ join has cost $U=V_i$.
At the cutoff sources in \eqref{eq:separation-costs}, the
logarithms of all loss factors in the fixed
raw packets, backgrounds, interpolations, and normalizations used
before the next join are $o(U)$.  Cutoff derivatives obey this
same rule.  At a current level-$i$ good join, a fixed
normalization can instead have logarithmic loss
$C\Re Z_i\leq C'\epsilon_iV_i$, which need not be $o(V_i)$.
Include these finitely many losses when choosing that join:
take a common sufficiently small fixed $\epsilon_i$ so that
$C'\epsilon_i$ uses at most half the available positive join
precision, increasing the finite transition budget when its
remainder is used.  The resulting join error is recorded as
$e^{-cV_i}$ with the reduced $c>0$.  Remaining later-rate,
interpolation, and fixed derivative losses have logarithmic
size $o(V_i)$.  All these choices precede variation of the
weight parameter below and are common for one controlled finite
request.  Already selected earlier joins remain fixed when
later requests add only later-rate normalizations there.
Local versions differ only by the stated errors.
\end{enumerate}
\end{definition}

The derivative strengthening used in Section~\ref{sec:calculus}
will impose these requirements also on each fixed finite collection
of independent derivatives.  No such strengthening is needed for
the separation theorem itself.

Pair the two supports by inserting zero packets and order their
terminal exponent vectors lexicographically, beginning with index
$1$.  Iterated left-finiteness implies that every nonempty set of
nonzero paired leaves has a first element: choose successively the
least coordinate admitting a nonzero descendant.  Zero formal data
always means that each terminal coefficient is zero as an analytic
germ, rather than just at $t_c$.

\begin{theorem}[Separation]\label{thm:separation}
Suppose the data of Definition~\ref{def:packet} satisfy its
requirements on a saturated flag.
If all paired terminal coefficients vanish, the actual function
is identically zero on the real flag locus sufficiently far in the
regime, locally in ordinary inputs.  Otherwise, let
$\lambda=(\lambda_1,\ldots,\lambda_m)$ be the first nonzero
paired exponent.  Its coefficients coincide as ordinary analytic
germs, say $c^+_\lambda=c^-_\lambda=c_\lambda$, and
\begin{equation}\label{eq:separation-leading}
 \exp\!\left(\sum_{i=1}^m\lambda_i Z_i\right)S^1
       \longrightarrow c_\lambda(t)
\end{equation}
locally uniformly on the real locus.

The same assertion holds for each selected leading prefix, with
all earlier whole subtrees zero.  In band $j$, use its raw packet
and normalize by the remaining selected exponents, evaluating all
these rates with $b_j$.  The coupled normalized determinations have
the same leading limit inside their selected contours and through
the used bands.  With
\eqref{eq:separation-background-fringe}, their leading estimates
extend through the owning outer fringe for any fixed required
normalization.  An entire zero subtree in band $p+1$, passed at
such a prefix, has size
\begin{equation}\label{eq:separation-zero-fringe}
 O(e^{-cV_p})
\end{equation}
on more vertical level-$p$ collars and on the subsequent portion
of that fringe, where its outer determination is in use.
All constants and contour choices concern fixed finite labels and
requests; none is asserted uniform over the whole packet tree.
\end{theorem}

To select a leading child at level $p$, we multiply its transition by
$e^{aZ_p}$, where $a$ is the selected exponent.  A child with an earlier
exponent $b<a$ then carries the growing factor $e^{(a-b)Z_p}$.
If that child's entire formal subtree is zero, decay to every fixed
order in its own, smaller scale is not yet an estimate that absorbs
this factor on the level-$p$ collar.  The essential analytic step is
to prove the stronger bound $O(e^{-cV_p})$ there.  Since
$\Re Z_p\asymp\epsilon_pV_p$ on that collar, a sufficiently small
fixed $\epsilon_p$ then absorbs all the finitely many earlier
multipliers needed for this transition.

\subsection{Coupling, holomorphic correction, and growth}

We begin with the fixed coupled prefix used in this argument.  Once
the exponents $\lambda_1,\ldots,\lambda_{p-1}$ have been isolated,
its function in band $j\leq p$ is
\begin{equation}\label{eq:separation-coupling}
 G=S^j_{\lambda_{<j}}
       \exp\!\left(\sum_{j\leq l<p}\lambda_lZ_{l,j}\right).
\end{equation}
The signs are understood on the two sides of the real ray, where
the two determinations share the same actual central function.
On a selected level-$j$ join, the adjacent formulas differ by
$O(e^{-cV_j})$.  We interpolate across its collar.  The logarithmic
cost of its derivatives is polynomial-logarithmic in the controlling
rates, and hence $o(V_j)$.  Its $\bar\partial$ source therefore
retains the same precision.

Each inner band $j<p$ stops at this selected good join.  Only the
outer band $p$ extends to its own side.  Consequently an inner
source satisfies
\begin{equation}\label{eq:separation-inner-cosine}
 \cos(\Im f_{j,j})\geq c/\Lambda_j.
\end{equation}
This lower bound will control a secondary cosine in
\eqref{eq:separation-costs} when its index is smaller than $p$.
Multiplication by any further fixed later normalization preserves
an old join error, with a smaller positive $c$, because the later
real phases are $o(V_j)$ on that collar.

We establish the zero-subtree bound by induction from the terminal
level toward the root.  Fix the functions and joins of a coupled prefix
through level $n$.  For an entirely zero subtree we first prove,
for each fixed $N$, the real estimate
\[
 |G(u)|\leq C_N e^{-NZ_n(u)}
 \qquad (u\geq u_N).
\]
The induction at deeper levels supplies this estimate by passing
only finitely many children for the chosen $N$.  Its constants and
starting point may depend on $N$; the already selected inner
functions and joins remain fixed.  The weighted maximum estimate
then removes those constants and gives the suppression at the
preceding transition.  At level $1$ there are no sources, and the
same estimate gives exact vanishing.  Once all zero subtrees can be
passed, a finite sequence of transitions isolates the first nonzero
leaf.  A comparison of its two lateral limits identifies its ordinary
coefficient.  The full induction is assembled at the end of the proof.

Figure~\ref{fig:coupled-bands} shows this domain and its join
sources.  The next two lemmas correct its later rates to holomorphic
functions and propagate the raw growth bound across the replaced
inner bands.
\begin{figure}[tbp]
\centering
\begin{tikzpicture}[
  x=1cm,y=1cm,
  every node/.style={font=\small},
  boundary/.style={semithick},
  annotation/.style={anchor=west,align=left,text width=3.6cm},
  direction/.style={-{Stealth[length=2mm]},thin}]
\fill[blue!4] (0,0) rectangle (8.4,1.05);
\fill[blue!8] (0,1.05) rectangle (8.4,2.25);
\fill[blue!4] (0,2.25) rectangle (8.4,4.55);
\fill[green!9] (0,3.35) rectangle (8.4,4.55);
\fill[orange!25] (0,0.95) rectangle (8.4,1.15);
\fill[orange!25] (0,2.15) rectangle (8.4,2.35);

\draw[boundary] (0,0)--(8.4,0);
\draw[boundary] (0,1.05)--(8.4,1.05);
\draw[boundary] (0,2.25)--(8.4,2.25);
\draw[boundary,dashed] (0,3.35)--(8.4,3.35);
\draw[boundary] (0,4.55)--(8.4,4.55);

\node at (4.2,0.5) {Used band $1$};
\node at (4.2,1.65) {Used band $2$};
\node at (4.2,2.82) {Outer band $3$};
\node at (4.2,3.92) {Level-$3$ outer fringe};

\node[annotation] at (8.6,0) {Real ray};
\node[annotation] at (8.6,1.05)
  {Selected $1$-good join\\collar cost $U=V_1$};
\node[annotation] at (8.6,2.25)
  {Selected $2$-good join\\collar cost $U=V_2$};
\node[annotation] at (8.6,3.35)
  {Available $3$-good contour\\not yet a join};
\node[annotation] at (8.6,4.55)
  {Outer side\\$\Re\widehat L_3=0$};

\draw[direction] (0,-0.4)--(8.4,-0.4)
  node[midway,below] {$u$ increases};
\node[anchor=west,font=\footnotesize] at (0,5.05)
  {Upper half of the $s$-domain; schematic, not to scale};
\end{tikzpicture}
\caption{Coupling through level $p=3$; the analogous lower half is
omitted.  Each inner band $j<3$ ends at its selected level-$j$ good
join, whose smoothing cost is $U=V_j$.  Only the outer band $3$
continues through its outer fringe to the level-$3$ side.  The dashed
$3$-good contour is available for the next transition and is not yet
a join of this coupling.  Shaded collars mark the two join sources;
interior cutoff sources are not displayed.}
\label{fig:coupled-bands}
\end{figure}


\begin{lemma}[Holomorphic proxies and Cauchy correction]
\label{lem:separation-proxies}
On a coupled level-$p$ domain, the rates of indices $l\geq p$
have holomorphic proxies $\widehat Z_l=Z_l+o(1)$, uniformly on a
slightly smaller domain.  Their logarithms have the prescribed
lifted determinations.  Where needed, $f_{p-1}$ also has an
absolute $o(1)$ holomorphic proxy.  The level-$p$ side can be
defined by $\Re\widehat L_p=0$, with
\[
 \widehat L_p=\widehat Z_p-A\sum_{l>p}\widehat Z_l
                 -\frac{\widehat Z_p}{(\log\widehat Z_p)^K}.
\]
All the angular comparisons and real-part domination estimates
of Section~\ref{sec:flags} remain valid on the resulting component
containing the real ray.

If $a(s)$ is a $\bar\partial$ source in a domain contained in a
fixed-width strip and $|a(s)|\leq M e^{-u}$ for $u\geq u_b$,
its Cauchy transform has norm at most $C M e^{-u_b}$, with $C$
depending only on the strip width, and tends uniformly to zero
as $u\to\infty$ on that fixed domain.  The norm bound is also
uniform over families with a common $M$ and strip width; in
particular it tends to zero when their left starts $u_b$ tend
to infinity.  These statements apply as well after multiplication
by a holomorphic gauge whenever its source has the displayed
envelope.
\end{lemma}

\begin{proof}
Here and below sources are extended by zero when forming their
area Cauchy transforms; the identity for $\bar\partial$ holds in
the interior.  For an evaluation point $s$, split the integral
against $1/(\pi(s-\zeta))$ into $|s-\zeta|<1$ and its complement.
The first kernel has bounded integral on the unit disk and the
second has modulus at most one.  Thus
\[
 \left|\frac1\pi\iint
        \frac{a(\zeta)}{s-\zeta}\,\dd A(\zeta)\right|
 \leq C\|a\|_\infty+C\|a\|_1
 \leq C M e^{-u_b}.
\]
Subtracting the transform makes a $C^1$ function holomorphic.
The small norm and decay on a fixed domain are distinct assertions.
For the latter, fix $R>u_b$ and split $a=a_{\leq R}+a_{>R}$
according to the real coordinate of its argument.  For $u>R+1$,
the first transform is bounded by
\[
 \frac{\|a_{\leq R}\|_1}{\pi(u-R)},
\]
uniformly in $v$.  The preceding near-kernel and far-kernel
estimate bounds the second transform by $CM e^{-R}$.
First let $u\to\infty$ with $R$ fixed, then let $R\to\infty$.
This proves the asserted absolute $o(1)$ on a fixed path without
asserting any pointwise exponential decay rate for the correction.
For varying anchors we instead use the uniform norm estimate
with their far-starting domains; its constant is independent of
the anchor when the source envelope and strip width are common.

We verify that the rate sources have the required small envelope.
A cutoff with leading index $\sigma(r)<p$, active on an inner
band $j$, has
\begin{equation}\label{eq:separation-early-logcost}
 \Re f_{r,j}\geq
 c\frac{Z_{\sigma(r)}(u)}{\operatorname{polylog}Z_{\sigma(r)}(u)}.
\end{equation}
Indeed, if $\Lambda_{\sigma(r)}\ll\Lambda_j$, its leading
cosine is bounded below; otherwise
$\Lambda_j=O(\operatorname{polylog}Z_{\sigma(r)})$ and
\eqref{eq:separation-inner-cosine} gives the assertion.
Later terms in its affine log relation are negligible by the
real-part comparisons of Lemma~\ref{lem:flags-contours}.
For a secondary cosine, the loss is bounded unless its rate
derivative is comparable with $\Lambda_j$.  In that case its
negative logarithm is $O(\log(2+\Lambda_k))$ and is absorbed by
\eqref{eq:separation-early-logcost}.  In particular the log-cost
dominates $f_l$ for $l>\sigma(r)$; for $\sigma(r)=p-1$ it
still dominates $f_{p-1}$, since
$f_{p-1}=\log Z_{\sigma(r)}$.

For $\sigma(r)\geq p$, positivity and the polynomial angular
buffer in \eqref{eq:separation-barrier} give $U\gg f_p$;
the more quantitative verification is included in
Lemma~\ref{lem:separation-source-ratios}.  Join costs also satisfy
$V_j\gg f_p$ for $j<p$.  All these statements hold uniformly
on the coupled domain; they imply in addition $U/e^u\to\infty$.
Differentiating a later rate through its triangular formula
introduces only a fixed finite product of later rates, of
logarithmic size $O(f_p)$, as well as the allowed background
losses.  These factors can therefore be absorbed into $e^{-cU}$.
The $\bar\partial$ derivatives of all rates $l\geq p$ are bounded
by arbitrarily small multiples of $e^{-u}$ far enough out.
The preceding integral constructs their absolute $o(1)$ proxies.

The continuous lifted logarithm of $Z_p$ and the estimate
$\widehat Z_p/Z_p=1+o(1)/Z_p$ give a continuous logarithm of
$\widehat Z_p$; since $\widehat Z_p$ is holomorphic and nonzero,
this logarithm is holomorphic.  Its difference from $f_p$ is
exponentially small in $f_p(u)$.  If $f_{p-1}$ is free it is
already holomorphic.  Otherwise its exact affine expression in
later rates and a bounded background admits the same correction.

Construct these corrections first on the piecewise domain with
extra room.  Near vertical, the inward rotation in $Z_p/f_p^K$
has positive real part of size at least
$c|Z_p|/f_p(u)^{K+1}$.  This tends to infinity faster than every
fixed power of $f_p(u)$ and overwhelms the absolute corrections.
The plane comparisons of Lemma~\ref{lem:flags-contours} show that
$\Re\widehat L_p$ reaches zero within the available room after
increasing $A$ and decreasing $K$.  Its positive component has
the asserted sides and comparisons.  Neither this argument nor
the maximum principles below requires a globally one-to-one
map $s\mapsto\log\widehat Z_p$.
\end{proof}

\begin{lemma}[Propagation of coupled growth]
\label{lem:separation-growth}
Suppose the first $p-1$ transitions have been isolated and coupled.
Their coupled function has the level-$p$ bound
\begin{equation}\label{eq:separation-coupled-growth}
 \log^+|G|\leq C\Re Z_p+CD_{p-1},
\end{equation}
with the exceptional term absent on the outer side and on
simultaneous safe columns.  If a level-$p$ term is isolated on a
good collar with exponentially small error and coupled to its
child, the corresponding level-$(p+1)$ function has the analogous
bound, after a fixed backstep.  At a terminal transition it is
bounded.  All old joins and source types are preserved.
\end{lemma}

\begin{proof}
The assertion starts with the raw estimates.  Write $n=p+1$.
On the chosen $p$ collar the new coupled function equals its raw
outer child up to $O(e^{-cV_p})$.  Divide it holomorphically by
$\exp(C'\widehat Z_n)$, choosing $C'$ to absorb the child's
$\Re Z_n$ growth.  Old sources, including derivatives of fixed
normalizations, remain exponentially small.  Subtract their
Cauchy transform.  The lateral boundary logarithm is then at
most $CD_p+O(1)$, while the old interior bound gives logarithmic
growth $O(\Re Z_p)$ plus a path constant.

Choose a common safe left column $u_b$ in a fixed backstep.
There the old bound costs $O(e^{f_p(u_b)})$.  Put
$X+iY=\log\widehat Z_p$.  On the good-contour domain the positive
harmonic functions
\begin{equation}\label{eq:separation-growth-barriers}
 \delta\Re(\widehat Z_p e^{a_0s}),\qquad
 C_0e^{f_p(u_b)}e^{-\kappa(X-f_p(u_b))}\cos(\kappa Y)
\end{equation}
are respectively an infinity barrier and a left boundary barrier.
Take $0<\kappa<1$.  The small $a_0>0$ is chosen relative to the
good-contour constant: $a_0|v|$ is smaller than half the angular
margin $\epsilon_p/\Lambda_p$, so the first real part is positive.
It dominates the growth at infinity for every fixed $\delta>0$.
Apply the logarithmic maximum principle on bounded truncated
components, then let the right boundary tend to infinity and
$\delta$ decrease to zero.  The remaining left contribution tends
to zero if a fixed backstep was chosen with
$f_p(u)>Bf_p(u_b)$ for a sufficiently large fixed $B$; such
backsteps exist by Lemma~\ref{lem:flags-path}.  Restoring the
divisor and the small Cauchy correction proves
\eqref{eq:separation-coupled-growth} with $p$ replaced by $n$.

The raw outer part already has the safe-column improvement.
To carry it through the replaced inner wedge, if $D_p$ is not
absorbed by the new rate bound, use the safe collar on the box
$|X-f_p(u_*)|\leq R$, where $R=C_1\log(2+D_p)$.
The end boundary loss $CD_p$ is majorized by a constant times
\[
 D_p\frac{\cosh(\kappa(X-f_p(u_*)))\cos(\kappa Y)}
               {\cosh(\kappa R)}.
\]
Choose the fixed $C_1$ so that this is bounded at $X=f_p(u_*)$.
The lateral collar has the improved bound by hypothesis; sources
remain small.  The box lies in the available range because
$R=o(f_p)$.  The maximum principle therefore removes $D_p$ on
the column.  The outer side improvement is raw.  At $p=m$ omit
the divisor and use bounded ordinary terminal coefficients;
the same proof gives boundedness.  A finite number of steps uses
only finitely many safe columns and labels, as required.
\end{proof}

\subsection{The source comparisons}

The growth bound is now available for a fixed coupled domain.  To
improve a zero subtree, we will multiply its function by
$\exp(H\widehat L_n-C\widehat Z_n)$, with $H>0$ and fixed $C$.
Write $x=\Re\widehat L_n$.  A source of precision $e^{-cU}$ then
has size at most $\exp(Hx-cU)$ after the fixed losses are absorbed.
Thus the ratio $U/x$ controls how large a weight $H$ the sources
permit.

For a target with value $x_a>0$, separate sources with
$x\geq\rho x_a$ from those with $x<\rho x_a$, where $\rho>0$
is fixed and small.  A lower bound for $U/x$ on the first set
pays the weight there; the second set costs at most
$\rho Hx_a$.  The following lemma gives the required lower bounds
in the two possible slope regimes.  It also proves a common
integrable envelope for a spare factor $e^{-cU/2}$.  This second
estimate makes the constant in the Cauchy correction bound independent
of $H$.

\begin{lemma}[Source ratios]\label{lem:separation-source-ratios}
In a coupled level-$n$ domain put $x=\Re\widehat L_n$.
For $n>1$ write
\[
 i=n-1,\qquad g=f_i-f_n,\qquad
 r_0=\frac{g'(u_a)}{\Lambda_n(u_a)},\qquad g_0=g(u_a).
\]
At a target on a sufficiently vertical level-$i$ collar, and
through its further fringe,
\begin{equation}\label{eq:separation-target-x}
 x_a\asymp\Re Z_n(s_a)
       \asymp e^{f_n(u_a)}\frac{r_0}{1+r_0}.
\end{equation}
For every fixed $\rho>0$ and every fixed upper bound on $r_0$,
every coupled source satisfying $x\geq\rho x_a$ obeys
\begin{equation}\label{eq:separation-ratio-small}
 \frac{U}{x}\geq c\frac{e^{g_0}}{g'(u_a)}.
\end{equation}
When $r_0$ is bounded below by a positive constant this also holds
for central targets with $x_a\asymp e^{f_n(u_a)}$.

For large $r_0$ and central targets with
$x_a\asymp e^{f_n(u_a)}$, one has instead
\begin{equation}\label{eq:separation-ratio-large}
 \frac{U}{x}\geq e^{c g_0}
 \quad\hbox{on sources with }x\geq\rho x_a.
\end{equation}
For a fixed coupled calculation, $U/e^u\to\infty$ and
$U/x\to\infty$ uniformly on its sources at infinity.
These comparisons are unchanged by the absolute proxy errors.
\end{lemma}

\begin{proof}
All estimates can first be made with the piecewise rates; their
buffer tends to infinity, so the $o(1)$ corrections have no
effect.  Write $t=f_n(u)$ and $\lambda=\Lambda_n(u)$, with a
subscript $0$ for the target.  From
Lemma~\ref{lem:flags-path},
\begin{equation}\label{eq:separation-rate-derivatives}
 c t\leq\lambda\leq t\operatorname{polylog}t,
 \qquad
 \left|\frac{\dd\log\lambda}{\dd t}\right|
       \leq\frac{\operatorname{polylog}t}{t}.
\end{equation}
At the target the angle of $f_n$ differs from that of $f_i$ by
the gain $g'/\Lambda_i$.  Since $g'\to\infty$, the fixed
good-contour margin is negligible in comparison.  Later real
phases are negligible against $\Re Z_n$.  This proves
\eqref{eq:separation-target-x}; the same is true if the level-$i$
angle is moved farther toward vertical.  Background jet errors
are swallowed by that gain, by
Lemma~\ref{lem:flags-background}.  In particular this gain is
larger than the polynomial buffer needed for the level-$n$ side.

\emph{Bounded target slopes and join sources.}
Suppose first that $r_0\leq R$ with $R$ fixed.
The source constraint and $x\leq Ce^t$ give
\begin{equation}\label{eq:separation-backward-t}
 t\geq t_0-O(1)-\log(1+1/r_0).
\end{equation}
Since $g'_0\geq cg_0\to\infty$,
$r_0\geq c/\lambda_0$, so the possible backward interval has
length $O(\log\lambda_0)=O(\log t_0)$.  On that interval
$\lambda\asymp\lambda_0$ by
\eqref{eq:separation-rate-derivatives}.  The gap estimates also give
\[
 \left|\frac{\dd\log g'}{\dd t}\right|
   =\left|\frac{g''}{g'\lambda}\right|
   \leq\frac{\operatorname{polylog}t_0}{t_0},
\]
because $g\leq g_0\leq C_R\lambda_0$ in a backward interval.
Consequently $g'\asymp g'_0$, and its integrated loss satisfies
\begin{equation}\label{eq:separation-gap-backstep}
 g(u)\geq g_0-O_R(1),
\end{equation}
since $r_0\log(1+1/r_0)$ is bounded on $(0,R]$.
At forward points,
\[
 \frac{\dd}{\dd u}\log\frac{e^g}{g'}
       =g'-\frac{g''}{g'}>0
\]
eventually, using $g'\geq cg$ and the polynomial-logarithmic
bound on $g''/g'$.  Thus in both directions
$e^g/g'\geq c_R e^{g_0}/g'_0$ whenever the source constraint
can matter.

A join at $j<n$ has $U=V_j$.  With $J=f_j-f_n$, the collar
comparisons give
\[
 x\leq C e^{f_n}\frac{1+J'}{\Lambda_j},\qquad
 \frac{U}{x}\geq c\frac{e^J}{J'}.
\]
Write $J=g+\Delta$, $\Delta=f_j-f_i\geq0$.  If $j=i$,
$\Delta=0$; otherwise $\Delta'\leq\Delta\operatorname{polylog}\Delta$.
Since $g'\to\infty$, exponential growth in $\Delta$ absorbs
the possible $\Delta'$ loss, giving
\[
 \frac{e^J}{J'}\geq c\frac{e^g}{g'}.
\]
This proves \eqref{eq:separation-ratio-small} for joins.

\emph{Cutoff sources at bounded target slopes.}
Consider a cutoff source, with earlier label $r<j\leq n$,
$F=f_r(u)$ and $\sigma=\sigma(r)$.  If $\sigma<n$, estimate
\eqref{eq:separation-early-logcost} applies on its own inner
band.  Because $\sigma\leq i$ and
$f_i\leq\log Z_\sigma$, its log-cost exceeds $f_i$ by an
unbounded factor.  The secondary cosine loss described there
does not change this conclusion.  It is therefore stronger
than both required ratios.

Suppose $\sigma\geq n$.  There are four estimates to check.
If $F\leq t^{3/2}$, the angular scale is at most
$C/\lambda$, and Taylor's formula together with the fixed
derivative bounds gives
\[
 \Re f_{r,j}=F+O\!\left(
              \frac{F\operatorname{polylog}F}{\lambda^2}
                         \right)=F+o(1).
\]
For a simple cost, or a mixed cost with $k\geq n$, division
by $x$ therefore gives
\begin{equation}\label{eq:separation-small-F}
 \frac{U}{x}\geq c e^{F-f_n}.
\end{equation}
For $k<n$, the source lies on an inner band $j\leq k$ ending
at its good join.  Its secondary cosine is at least
$c/\Lambda_k$, and the angular comparison yields
\begin{equation}\label{eq:separation-cosine-ratio}
 \frac{\cos(\Im f_{k,j})}{\cos(\Im f_{n,j})}
       \geq\frac{c}{1+\Lambda_k-\Lambda_n}.
\end{equation}
Since $r<j\leq k\leq i$, we have $F\geq f_k$.
Set $\Delta=f_k-f_i\geq0$.  The resulting ratio is at least
\[
 c\frac{e^{g+\Delta}}{1+g'+\Delta'}
       \geq c\frac{e^g}{g'}.
\]
Together with the backward and forward comparisons already
proved, this settles the small-$F$ case.

If $t^{3/2}<F\leq t^3$, the leading affine relation gives
$F\asymp Z_\sigma$ and
$f_\sigma=\log F+O(1)=O(\log t)$.  Hence
$\Lambda_\sigma/\lambda\to0$, the angle of its leading rate
is small, and
\begin{equation}\label{eq:separation-middle-F}
 \Re f_{r,j}\geq cF.
\end{equation}
If $F>t^3$ and $\sigma>n$, the leading cosine is at least
$c/\lambda$; later affine terms are negligible.  Thus
\begin{equation}\label{eq:separation-large-F}
 \Re f_{r,j}\geq cF/\lambda.
\end{equation}
Finally, if $\sigma=n$, the large $A$ in the side equation
absorbs the later affine terms and leaves the buffer:
\begin{equation}\label{eq:separation-sigma-n}
 \Re f_{r,j}\geq cx+
                    c\frac{e^t}{t^{K+1}}.
\end{equation}
This also holds on the relevant inner contours.

For bounded target slopes,
$g_0\leq C_R\lambda_0$.  At backward sources the possible
drop in $t$ is only $O(\log t_0)$; at forward sources
$g_0\leq C_R\lambda_0\leq t\operatorname{polylog}t$.
Each of \eqref{eq:separation-middle-F}--\eqref{eq:separation-sigma-n}
therefore dominates $t+g_0+\log(2+F)$.  In
\eqref{eq:separation-large-F} use
$F/\lambda\geq t^2/\operatorname{polylog}t$ and
$\log F\leq t+O(1)$ when $\sigma>n$.
The negative logarithm of a secondary cosine is either bounded
or $O(\log(2+F))$: if $k<n$, use its inner-band lower bound;
if $k\geq n$, use the side's polynomial buffer and ordered
angles.  All those losses and division by $x\leq Ce^t$ are
absorbed.  This proves \eqref{eq:separation-ratio-small} for
all cutoff costs, uniformly even when $r_0\to0$.

\emph{Large target slopes.}
Now take a large target slope and a central target.  The source
constraint implies $t\geq t_0-O(1)$.  The adjacent-gap bound
$g'/g\leq C\lambda$ implies
\begin{equation}\label{eq:separation-large-backstep}
 g(u)\geq c g_0
 \quad\hbox{whenever }t\geq t_0-O(1).
\end{equation}
The join calculation and the small-$F$ calculation above give
$U/x\geq e^{c g_0}$: their derivative denominators are only
polynomial-logarithmic in the corresponding diverging gaps.
For $\sigma=n$, the target relation gives
$g_0\leq f_i(u_a)\leq f_r(u_a)\leq Ce^{t_0}$.
Thus \eqref{eq:separation-sigma-n} and the source constraint
$x\geq\rho x_a\asymp e^{t_0}$ leave a fixed positive multiple
of $g_0$ in the log-cost.  The buffer term absorbs $t$ and the
secondary cosine losses, giving the same ratio.

It remains to justify the extra angular loss when $F>t^{3/2}$
and $\sigma>n$.  Put
\[
 J=f_n-f_\sigma,\quad
 d=\frac{J'}{\lambda}=1-\frac{\Lambda_\sigma}{\lambda},\quad
 \eta=\cos(\Im f_{n,j}),\quad q=\min(1,\eta+d).
\]
The leading real contribution satisfies
\begin{equation}\label{eq:separation-angular-F}
 \Re f_{r,j}\geq cF(\eta+d).
\end{equation}
Moreover $d\geq c/\lambda$, and so
$L=\log(1/q)=O(\log t)$.  The gap derivative bounds give
\begin{equation}\label{eq:separation-logd}
 \left|\frac{\dd\log d}{\dd t}\right|
 =\left|\frac{J''}{J'\lambda}
                    -\frac{f_n''}{\lambda^2}\right|
       \leq\frac{\operatorname{polylog}t}{t}.
\end{equation}
Decrease $t$ to $t-L$ and denote the corresponding real position
by $u'$.  The quantity $d$ changes by a bounded factor (in fact
by $1+o(1)$).  Since $q\geq d$,
$dL\leq d\log(1/d)$ is bounded.  Integrating
$\dd f_\sigma/\dd t=1-d$ now gives
\begin{equation}\label{eq:separation-angular-backshift}
 Z_\sigma(u')
  =Z_\sigma(u)
     \exp\!\left(-L+\int_{t-L}^{t}d(\tau)\,\dd\tau\right)
  \asymp qZ_\sigma(u).
\end{equation}
Because $x\leq Ce^t\eta\leq Ce^tq$ and
$x\geq\rho x_a$, one has $t-L\geq t_0-O(1)$.
The indices $r<j\leq n=i+1$ imply $r\leq i$; hence
\[
 g(u')\leq f_i(u')\leq f_r(u')\asymp Z_\sigma(u').
\]
Combining this with \eqref{eq:separation-large-backstep} and
\eqref{eq:separation-angular-backshift} proves
$F(\eta+d)\geq c g_0$.
Independently, \eqref{eq:separation-middle-F} or
\eqref{eq:separation-large-F} gives
$\Re f_{r,j}\gg t+\log(2+F)$.
These two lower bounds, each used with a sufficiently small
fixed fraction of the log-cost, absorb division by $x$ and all
secondary cosine losses, leaving $\log(U/x)\geq c g_0$.
This proves \eqref{eq:separation-ratio-large}.

\emph{The common integrable envelope.}
The same estimates without a target show the asserted uniform
growth at infinity.  For joins,
$U/x\geq ce^{f_j-f_n}/(f_j-f_n)'\to\infty$.
For cutoffs with small $F$, the bound is of this form with a
larger gap; the other cases have a log-cost dominating $f_n$
or its relevant angularly diminished value.  All costs dominate
$e^u$ uniformly over the controlled finite-request families in
Definition~\ref{def:packet}.  More explicitly, on a join
\[
 \log U=f_j-\log\Lambda_j\geq\tfrac12 f_j
                         \geq\tfrac12 e^u
\]
eventually.  For a cutoff, the small-$F$ estimate leaves
$F-O(\log(2+F))\geq c f_n$ after the secondary cosine loss;
the intermediate and large-$F$ estimates and the buffer in
\eqref{eq:separation-sigma-n} leave at least $c f_n$ as well.
The early-index estimate is stronger.  Thus, after decreasing
a fixed $c>0$ for the finite source list,
\begin{equation}\label{eq:separation-uniform-source-envelope}
 U\geq\exp(c e^u).
\end{equation}
Every threshold here depends only on the finite quantitative
controls specified in the definition.  In particular, for any
fixed source precision constant $c_*>0$, the residual factor
$e^{-c_*U}$ has a common envelope $M e^{-u}$ on the full
forward domains, with $M$ independent of their anchors.
The stipulated uniform $o(U)$ absorption permits the same
conclusion after all fixed losses.  These constants are also
independent of any subsequently chosen weight parameter.
\end{proof}

\subsection{The zero-subtree estimate}

\begin{lemma}[Weighted maximum estimate]\label{lem:separation-weighted}
Let $G$ be coupled through level $n$ with the growth bound of
Lemma~\ref{lem:separation-growth}.  Suppose that on the real ray
\begin{equation}\label{eq:separation-preliminary}
 |G(u)|\leq C_N e^{-NZ_n(u)}
 \quad(u\geq u_N),\qquad\hbox{for every fixed }N>0.
\end{equation}
The constants $C_N,u_N$ can depend arbitrarily on $N$.
Fix the coupled functions, its joins, a safe left column, and
the full domain $x=\Re\widehat L_n>0$.  There are fixed
$C,C_0,c>0$ such that for every $H>0$ and every point of that
domain,
\begin{equation}\label{eq:separation-weighted}
 \log|G|+Hx-C\Re\widehat Z_n
 \leq C_0+\max\left\{
 H\sup_{\rm left}x,
 \sup_{\rm sources}(Hx-cU/2),\ 0\right\}.
\end{equation}
In particular $C_0$ is independent of $H$, $C_N$, and $u_N$.
When $n=1$ there are no sources, and
\eqref{eq:separation-preliminary} implies exact vanishing.
\end{lemma}

\begin{proof}
\emph{Correcting the weighted sources.}
Choose $C$ greater than the fixed coefficient in the coupled
growth estimate on safe columns and the side.  Multiply $G$ by
the holomorphic gauge
\[
 \exp(H\widehat L_n-C\widehat Z_n).
\]
Its side modulus is bounded independently of $H$ because $x=0$;
its left boundary logarithm is at most
$H\sup_{\rm left}x+O(1)$.  With constants decreased to absorb
the fixed losses, its sources have modulus at most a finite sum
of $\exp(Hx-cU)$.  Set
\[
 K_H=\max\{0,\sup_{\rm sources}(Hx-cU/2)\}.
\]
The source-ratio lemma makes $K_H$ finite for every fixed $H$.
Crucially, keep half the negative cost:
\begin{equation}\label{eq:separation-spare-cost}
 e^{Hx-cU}\leq e^{K_H}e^{-cU/2}.
\end{equation}
The last factor is bounded by a fixed multiple of $e^{-u}$ and
is integrable on the bounded-width domain, because
$U/e^u\to\infty$.  Lemma~\ref{lem:separation-proxies} thus
gives a Cauchy correction of norm at most $C_1e^{K_H}$,
where $C_1$ is independent of $H$.  Let $F_H$ be the resulting
holomorphic function.
By \eqref{eq:separation-uniform-source-envelope}, $C_1$ is
uniform for the controlled finite-request families as well;
the bound uses their full forward domains, not just a window
about the target.  Their possibly different return times to
smaller ordinary neighborhoods do not enter this constant.

\emph{Boundedness for each fixed weight.}
For each fixed $H$, preliminary real decay makes the weighted
function bounded on the real ray: choose $N$ larger than the
fixed multiple of $H+C$ needed there, and include the finite
interval before $u_N$ in its bound.  This bound can depend on
$H$ in an uncontrolled way; we do not use it in the final
estimate.  To see that $F_H$ is bounded on each half-domain,
write $X+iY=\log\widehat Z_n$.  Its logarithmic growth is
$O_H(e^X)$ plus a path constant.  On the upper half use the
positive harmonic barrier
\[
 \delta e^{aX}\cos(a(Y-\pi/4)),\qquad 1<a<2,
\]
and on the lower half use
$\delta e^{aX}\cos(a(Y+\pi/4))$.
The side angles tend to $\pm\pi/2$, so a sufficiently far
start makes both barriers positive on their respective
half-domains.  They dominate the infinity growth.  The maximum
principle on finite truncations and then $\delta\downarrow0$
proves boundedness on each half, hence on the full domain.

\emph{Removing the preliminary constants.}
Now apply the bounded maximum principle to that \emph{full}
domain.  The real ray is no longer a boundary.  One can justify
the infinite-domain principle by
$\delta e^{a'X}\cos(a'Y)$, $0<a'<1$, before letting
$\delta\downarrow0$.  Only the side and the fixed safe left
column remain.  After restoring the Cauchy correction their
bounds give \eqref{eq:separation-weighted}, with an additive
constant depending on $C_1$ and the fixed data alone.  This
explains why the arbitrary constants in
\eqref{eq:separation-preliminary} disappear.

For $n=1$, choose a fixed backstep so that the target's positive
$x_a$ is larger than $\sup_{\rm left}x$.
There is no source term, and
\eqref{eq:separation-weighted} gives
\[
 \log|G(s_a)|\leq C_0+C\Re\widehat Z_1(s_a)
              -H(x_a-\sup_{\rm left}x).
\]
Letting $H\to\infty$ proves $G(s_a)=0$.  Every sufficiently
far real target permits this construction, giving the asserted
identity on the real locus.
\end{proof}

\begin{lemma}[Upgrade at the preceding transition]
\label{lem:separation-upgrade}
For $n>1$, under the hypotheses of
Lemma~\ref{lem:separation-weighted}, the outer packet of $G$
satisfies $|G|\leq Ce^{-cV_{n-1}}$ on sufficiently vertical
level-$(n-1)$ collars and their further fringe.
\end{lemma}

\begin{proof}
Put $i=n-1$ and use the notation of the source-ratio lemma.
A fixed backstep can make
\begin{equation}\label{eq:separation-left-small}
 \sup_{\rm left}x\leq\rho x_a
\end{equation}
for any chosen small fixed $\rho>0$.
Indeed, $f_n$ on that column can be made a fixed factor smaller
than $f_n(u_a)$, whereas
$x_a\geq c e^{f_n(u_a)}/\Lambda_n(u_a)$ at the target.
Safe columns exist in the needed fixed intervals.

\emph{Bounded slopes.}
If $r_0$ is bounded, choose
\[
 H=c_1\frac{e^{g_0}}{g'_0}
\]
with $c_1>0$ sufficiently small.  On sources with
$x\geq\rho x_a$, \eqref{eq:separation-ratio-small} gives
$Hx-cU/2\leq0$; on the remaining sources it is at most
$\rho Hx_a$.  Thus the maximum on the right of
\eqref{eq:separation-weighted} is at most $\rho Hx_a$.
At the target $\Re\widehat Z_n\asymp x_a$ and $H\to\infty$,
so the subtracted fixed growth term is negligible.  Finally,
\begin{equation}\label{eq:separation-bounded-slope-scale}
 Hx_a\asymp
 \frac{e^{g_0}}{g'_0}e^{f_n(u_a)}
        \frac{g'_0/\Lambda_n(u_a)}{1+g'_0/\Lambda_n(u_a)}
 =\frac{e^{f_i(u_a)}}{\Lambda_i(u_a)}=V_i(u_a).
\end{equation}
This proves the result for bounded slopes, uniformly down to
$r_0=0$ as a limiting regime.  When $r_0$ stays in a fixed
compact subinterval of $(0,\infty)$ the same proof applies
to central level-$n$ targets.

\emph{Large slopes: decay in the central region.}
For large slopes, first use central level-$n$ targets.
Take $H=\exp(c_2g_0)$ with $c_2$ smaller than the constant
in \eqref{eq:separation-ratio-large}.  Exactly the preceding
argument gives
\begin{equation}\label{eq:separation-central-large}
 |G(s_a)|\leq
       \exp\{-\exp(c_3 f_i(u_a))\}
\end{equation}
for some fixed $c_3>0$, since
$Hx_a\geq c\exp(f_n(u_a)+c_2g_0)$ and
$f_i=f_n+g_0$.  This estimate is weaker than the desired one
at the preceding transition, so a second maximum principle is
needed.

\emph{Large slopes: transfer to the preceding fringe.}
Use the holomorphic coordinate function
$X+iY=\widehat f_i$.  Work first in the local tube
$|v|<2W/\Lambda_i(u)$, with fixed slack at its real end columns,
and take the component meeting the real ray of the inverse
image, within this tube, of the rectangle
\[
 X_0-1<X<B_1X_0,\qquad |Y|<W,
 \qquad X_0=f_i(u_a),
\]
where $W$ is a large fixed number.  Stop at a right real column
if $\Lambda_i/\Lambda_n$ first falls to a fixed
$M_0\gg W$.  For a sufficiently large starting ratio, its
short left extension retains a large ratio by the derivative
bounds.  On the tube's two angular walls the path estimates and
the absolute proxy correction give
$Y=\pm2W+o(1)$ and $X=f_i(u)+o(1)$.
The selected component, where $|Y|<W$, cannot meet those walls.
There are therefore no additional angular boundary pieces to
estimate.  All its points remain in this local tube and are
central for level $n$: their $f_n$ angle is at most
$O(W/M_0)$.
The level-$i$ target, with $|Y|\sim\pi/2$, is inside it.

Every source encountered in this rectangle has
\begin{equation}\label{eq:separation-rectangle-cost}
 U\geq c\frac{e^X}{\Lambda_i}.
\end{equation}
For joins this is the gap calculation above.  A cutoff with
$\sigma<n$ has the stronger bound
\eqref{eq:separation-early-logcost}.  If $\sigma\geq n$,
its leading angle is small.  With $F=f_r$ and $T=f_i$, the
case $F\leq T^{3/2}$ gives $\Re f_{r,j}=F+o(1)$ by the
Taylor calculation on the scale $W/\Lambda_i$.  Since
$r\leq i$, this pays $e^X$; a secondary cosine with $k\leq i$
costs at worst $\Lambda_k$, and
$\Lambda_k/\Lambda_i$ is absorbed by $e^{F-f_i}$ using
the same gap calculation as
\eqref{eq:separation-cosine-ratio}.  If $F>T^{3/2}$, use
that every phase of index at least $n$ has size
$O(W/M_0)$ throughout this rectangle.  The leading cosine
is bounded below, later affine terms remain negligible,
and $\Re f_{r,j}\geq cF$.  This exceeds
$T+\log(2+F)$ by an unbounded factor, paying every such
secondary loss.  This proves
\eqref{eq:separation-rectangle-cost}.

Take $\kappa=\pi/(2W)<c_3$, and let $w$ be the positive
harmonic function
\begin{equation}\label{eq:separation-strip-weight}
 w(X,Y)=c_4\frac{e^{X_0}}{\Lambda_i(u_a)}
      \sinh(\kappa(X-X_0+1))\cos(\kappa Y).
\end{equation}
It is the real part of a holomorphic function of
$\widehat f_i$ and vanishes on the left and horizontal sides.
The rate derivative bound implies
\begin{equation}\label{eq:separation-lambda-strip}
 \left|\frac{\dd\log\Lambda_i}{\dd X}\right|
      \leq\frac{\operatorname{polylog}X}{X}=o(1).
\end{equation}
Consequently, throughout the rectangle,
\[
 \frac{w(X,Y)}{e^X/\Lambda_i(X)}
 \leq Cc_4
       \exp(-(1-\kappa-o(1))(X-X_0))
\]
apart from the harmless one-unit left extension.  Choosing $c_4$
small makes the weight at most a fixed small fraction of every
source cost in \eqref{eq:separation-rectangle-cost}.

On a full right edge $X=B_1X_0$, the logarithm of the size of
$w$ is at most
$[1+\kappa(B_1-1)]X_0+O(\log\Lambda_i(u_a))$.
Choose
\begin{equation}\label{eq:separation-right-factor}
 B_1>\frac{1-\kappa}{c_3-\kappa}.
\end{equation}
Then \eqref{eq:separation-central-large} pays the whole
weight there.  If the domain was stopped at the bounded-ratio
column, the bounded-slope result
\eqref{eq:separation-bounded-slope-scale} pays it instead,
after reducing $c_4$ using
\eqref{eq:separation-lambda-strip}.  The other edges have
zero weight, and \eqref{eq:separation-central-large} bounds
$G$ there.

Multiply $G$ by the holomorphic gauge with real logarithm $w$.
Its sources retain an integrable negative half-cost, so their
Cauchy correction is bounded independently of the target.
The maximum principle on the bounded component therefore
bounds the weighted function by a fixed constant.  At the
target, the cosine in \eqref{eq:separation-strip-weight} is
bounded below and $\sinh\kappa>0$, so
$w(s_a)\geq cV_i(u_a)$.  This proves the desired
$e^{-cV_i}$ bound.  The argument uses only positive harmonic
functions pulled back by $\widehat f_i$ and works without
global coordinate injectivity.  It applies unchanged to the
farther level-$i$ fringe.  Nominal pointwise cutoff variants
differ there by errors with the same lower cost, so satisfy
the same bound.
\end{proof}

\subsection{Finite isolation and terminal comparison}



\begin{lemma}[One finite transition]\label{lem:separation-isolate}
Suppose a prefix is coupled through level $p$.  Fix a child
exponent $a$, and assume that every earlier child has an entire
zero subtree and satisfies the upgrade of
Lemma~\ref{lem:separation-upgrade} at level $p$.
Then a sufficiently vertical fixed good collar can be chosen so
that
\begin{equation}\label{eq:separation-isolate}
 e^{aZ_{p,p}}S^p_{\lambda_{<p}}
     =S^{p+1}_{\lambda_{<p},a}+O(e^{-cV_p})
\end{equation}
on both sides, with zero insertion if necessary.  It can be
coupled to that child while preserving all preceding joins.
\end{lemma}

\begin{proof}
Only finitely many children precede $a$ by left-finiteness.
For an earlier exponent $b<a$, its outer packet is
$O(e^{-c_bV_p})$ beyond its previously selected join.
On a more vertical good collar,
$\Re Z_p\asymp\epsilon_pV_p$.  Choose $\epsilon_p>0$
small enough that every finite multiplier
$\exp((a-b)\Re Z_p)$ is absorbed by half of the corresponding
$c_bV_p$.  This is a finite restriction.  The collar may also
be chosen beyond every earlier child's selected join and within
the allowed inward extension of the chosen child.

Take a transition budget larger than $a$ and use
\eqref{eq:separation-transition}.  For every included later
exponent $b>a$, its positive gap $b-a$ suppresses its raw
packet: on the $p$ collar
\[
 \Re Z_{p+1}=o(V_p),\qquad D_p=o(V_p),
\]
and there are only finitely many such terms.  The first estimate
is the contour hierarchy.  For the second, any uncapped
coefficient contributes at most $O(\log\Lambda_p)$ on a fixed
backstep, and $\log\Lambda_p=o(V_p)$; the same conclusion
holds along the ray.  The tail error is also
$O(e^{-cV_p})$ after increasing the budget if necessary.
These observations give \eqref{eq:separation-isolate}.

In the inner bands multiply the existing coupled function by
$\exp(a Z_{p,j})$, then switch on this collar to the selected
outer child.  On an old $j$ join, $j<p$, the new rate has real
part $o(V_j)$, so old seam errors persist.  Background changes
alter this fixed multiplier by a relatively exponentially small
amount there.  The source hypotheses handle all derivatives of
the interpolation and normalization.  The growth bound for the
new coupling follows from Lemma~\ref{lem:separation-growth}.
No new transition at an earlier inner band has been requested.
\end{proof}

\begin{lemma}[Comparison of terminal coefficients]
\label{lem:separation-terminal}
Suppose a complete exponent vector has been isolated on the two
sides and coupled, with all earlier subtrees passed as in
Lemma~\ref{lem:separation-isolate}.  Its two terminal ordinary
coefficients are the same analytic germ.  The fully normalized
coupled function tends to this germ on the real ray and through
the corresponding inner wedge and used bands.
\end{lemma}

\begin{proof}
By terminal growth propagation the normalized function is
bounded on the whole inner domain after a fixed backstep.
For a target family with real anchors $u_a\to\infty$, choose
safe left columns $u_b$ in fixed backstep intervals.  The finite
raw, background, and source data used by this terminal coupling
have common bounds on their full forward domains, by the
finite-request hypothesis.  In particular
\eqref{eq:separation-uniform-source-envelope} supplies a common
$M e^{-u}$ envelope for their corrected sources.  Their Cauchy
corrections therefore have norm at most $CM e^{-u_b}\to0$,
uniformly over the target family.  On a single fixed path the
compact-left and tail argument in
Lemma~\ref{lem:separation-proxies} gives the corresponding
vanishing at infinity.  Denote the corrected functions by $\widetilde G$; they are
holomorphic and uniformly bounded.
On the two outer collars their values approach, respectively,
$c_+(t)$ and $c_-(t)$, by the terminal transition.

We give the argument that these two boundary limits cannot
differ.  Fix an ordinary value $t_*$ in a sufficiently small
real neighborhood and take a family of tests approaching it
on every fixed $u$ window about the varying anchor.  Such tests
can return to $t_c$ at infinity: for example, on the needed
right half-strip use
\[
 t(s)=t_c+(t_*-t_c)e^{-\delta_a(s-u_a)},
 \qquad\delta_a\downarrow0,
\]
with the fixed backward range and ordinary neighborhood chosen
with slack.  Its fixed positive derivatives tend to zero and
it returns to $t_c$ on every chosen path.  Taking
$0<\delta_a\leq\delta_0$ gives common finite-jet bounds on
the entire forward domains, and $\delta_a\to0$ gives a
common small-jet modulus and convergence to $t_*$ on each
fixed target window.  The time of return to a smaller box may
diverge; that time is used only in the preliminary real
estimates, whose constants have already been removed by
Lemma~\ref{lem:separation-weighted}.

Write $X+iY=\log\widehat Z_m$.  For a fixed small $\eta>0$,
the boundary values on a window about the target are within
$\eta$ of the corresponding constants, once the anchor is
sufficiently far.  The $X$ distances from the target to both
ends of a fixed $u$ window tend to infinity by
Lemma~\ref{lem:flags-path}.  On that window apply the logarithmic
maximum principle to the difference from $c_+(t_*)$.
Its upper boundary is bounded by $\log\eta$, its lower
boundary by a fixed constant $\log M$, and its two ends by
the bounded growth estimate.  An affine function of $Y$
interpolates the lateral logarithmic bounds.  The actual angles
tend to $\pm\pi/2$; adding arbitrarily small relative angular
slack and decaying cosine barriers at the two ends yields, at
the real target,
\[
 \limsup\log|\widetilde G-c_+(t_*)|
       \leq\tfrac12\log\eta+\tfrac12\log M.
\]
For completeness, the end barriers may be taken as fixed
multiples of
$e^{-\kappa(X-X_-)}\cos(\kappa Y)$ and
$e^{-\kappa(X_+-X)}\cos(\kappa Y)$ with $0<\kappa<1$.
Their values at the target tend to zero.  First let the target
go to infinity, then let the angular slack and $\eta$ go to
zero.  The real-ray limit is $c_+(t_*)$.  Applying the same
argument from the lower boundary gives the limit $c_-(t_*)$.
Thus these two values coincide for every such real $t_*$.
The identity theorem for analytic germs makes the coefficients
equal.

With their common value on both lateral boundaries, apply the
same maximum principle to their difference from $\widetilde G$.
The lateral bound is now $\eta$ on both sides, so the estimate
holds uniformly across the entire inner wedge, including
targets approaching either boundary.  Restoring the vanishing
Cauchy correction changes nothing.  Uniformity on smaller
ordinary neighborhoods follows either from these estimates
with common finite slack or by applying the construction to
a putative sequence of violating ordinary inputs.  The same
argument uses the corresponding coupled determinations in the
already used bands, and proves their asserted convergence.
\end{proof}

\begin{proof}[Proof of Theorem~\ref{thm:separation}]
We describe the induction explicitly to keep its finite and
infinite requests separate.  A coupled prefix includes fixed
functions on the earlier inner bands, fixed selected joins,
the growth bounds, and the finite list of their source types.
An assertion about that prefix may request more accuracy
\emph{at infinity on its existing path}; it may not change
those already fixed functions or joins.

Induct downward on the level $n$.  The assertion to prove is:
for every coupled prefix whose entire level-$n$ subtree has
zero terminal data, the preliminary estimates
\eqref{eq:separation-preliminary} hold for all fixed $N$;
they imply the preceding-transition upgrade when $n>1$ and
exact vanishing when $n=1$.

At $n=m$, all ordinary child coefficients are zero.  For each
fixed budget the transition formula therefore gives arbitrary
exponential decay on a sufficiently vertical $m$ collar.
Equivalently insert a zero terminal child at any exponent
$a>N$, choosing the remainder budget greater than $a$.
The proof of growth propagation, applied after multiplication
by $e^{aZ_m}$, gives a bounded function in the inner wedge
after a fixed backstep.  On the real ray this is
$O(e^{-aZ_m})$, and hence implies
\eqref{eq:separation-preliminary} for the requested $N$.
No uniformity over $N$ has been used.  The weighted lemma and
the upgrade lemma now establish the induction assertion at
level $m$.

Suppose it is known strictly below level $n$ and consider an
entirely zero subtree coupled through level $n$.
Fix one finite requested real order $N$.  Choose an exponent
$a>N$ at this coordinate and insert a zero child there if
necessary, enlarging the transition budget beyond it.  There
are finitely many children earlier than $a$.  Pass them in
increasing order.  To pass one child, first isolate it on its
own sufficiently vertical good contour using
Lemma~\ref{lem:separation-isolate} and the upgrades already
obtained for its earlier siblings.  Its coupled subtree at
level $n+1$ is zero, so the induction hypothesis below $n$
gives its preliminary real estimates and its
$e^{-cV_n}$ upgrade.  This permits the next sibling to be
isolated.  Only finitely many siblings are needed before
reaching the inserted zero child $a$.

After isolating that child, the growth propagation lemma gives
on the real ray
\[
 |G(u)|\leq C\exp(-aZ_n(u)+CZ_{n+1}(u)+CD_n)
\]
when $n<m$, and the corresponding bound without the last
rate at $n=m$.  Since $Z_{n+1}=o(Z_n)$ and $D_n$ is a
constant on this path, this is the requested
$O(e^{-NZ_n})$ estimate sufficiently far out.  Alternatively
one may continue the zero child down to a terminal zero;
both constructions use only finitely many transitions for
this request.  We have proved all fixed preliminary orders
at level $n$, and can now apply
Lemma~\ref{lem:separation-weighted} and
Lemma~\ref{lem:separation-upgrade}.

There is no circular use of a stronger estimate at the current
level.  Preliminary orders at level $n$ invoke upgrades only
at level $n+1$; the height decreases each time.  Finite breadth
can grow with $N$, but the height is at most $m$.
Moreover all joins strictly preceding the child under
consideration remain fixed.  Higher requests need new labels
only at that child or deeper, and any new fixed later
multiplier costs $o(V_j)$ on an old $j$ join.  Thus the full
domain in the weighted lemma has fixed source estimates
independent of its multiplier $H$.  The preliminary real
estimates may start arbitrarily farther out as $H$ varies;
their constants disappear in
\eqref{eq:separation-weighted}.  This completes the downward
induction and, at $n=1$, proves the exact-zero assertion.

If the paired array is nonzero, choose its first nonzero
exponent $\lambda$.  At the first coordinate pass only the
finitely many earlier subtrees, all of which are zero by
definition of $\lambda$.  The established upgrades and
Lemma~\ref{lem:separation-isolate} isolate $\lambda_1$.
Repeat at its selected child, then at the next one.  This
finite process reaches the paired ordinary coefficients at
$\lambda$.  Lemma~\ref{lem:separation-terminal} proves their
equality and the real leading limit
\eqref{eq:separation-leading}.  It also proves the corresponding
limits for the coupled normalized determinations inside their
selected contours and along the bands already used.

It remains to justify the owning-fringe extension in the
statement.  Beyond a selected leading join the outer packet
is the one whose convergence was just proved.  In the finite
transition \eqref{eq:separation-transition}, every earlier
term still has its $e^{-cV_j}$ suppression on that more
vertical portion, and every later term has a strictly positive
current exponent.  Choose the fringe constant sufficiently
small for the finitely many earlier multipliers and choose
$A$ large enough that $\Re Z_j$ dominates the finite later
real-phase losses.  The transition remainder can be requested
to any fixed needed depth.  Thus the same normalized leading
estimate holds throughout that fringe.

In this last comparison rates in the two backgrounds must
also be compared.  Estimate
\eqref{eq:separation-background-fringe} is sufficient: triangular
differentiation of a fixed later rate $Z_k$, $k>j$, has
logarithmic sensitivity at most $C f_{j+1}$ in these variables,
whereas $\Re Z_j$ dominates every fixed power of $f_j$ on the
buffered side.  Hence the absolute change in each such later
rate is $o(1)$ after any fixed needed normalization.  The
extracted current factor is always kept at its own inner
background; it is not reevaluated at the next background.
On selected good joins the stronger $e^{-cV_j}$ comparison
already gave the same conclusion.  This proves the stated
fringe extension and finishes the theorem.
\end{proof}

\subsection{Uniformity for the differentiated calculus}

The preceding proof also supplies the uniform estimates needed when
ordinary tests vary with their anchors.  We collect their precise scope
here for the differentiated calculus in Section~\ref{sec:calculus}.

\begin{lemma}[Uniformity for a controlled finite request]
\label{lem:separation-family-uniformity}
Fix finitely many packet labels, transition budgets, coupled
joins, and normalizations, and a family of tests with the common
quantitative controls in Definition~\ref{def:packet}.  The fixed
labels used on the coupled domain must stay in their common
ordinary ranges throughout that domain.  Deeper labels used
only for preliminary real estimates may have smaller ranges
entered at a time depending on the member of the family.

The proxy, growth, source-ratio, and preceding-transition
estimates above have constants uniform in this family.  If
each member has the preliminary estimates
\eqref{eq:separation-preliminary}, no uniformity of their
$C_N,u_N$ or of their return times is needed for the weighted
estimate or its upgrades.  Terminal comparison is uniform at
receding anchors when the tests have a common slow-variation
modulus on each fixed window about those anchors.  The same
statements apply to any fixed finite collection of derivative
packet trees satisfying the differentiated hypotheses.
\end{lemma}

\begin{proof}
First choose common sufficiently vertical collars and residual
positive join precisions after the finite current-level losses
have been absorbed as specified in Definition~\ref{def:packet}.
Take the largest of the finitely many raw-growth constants,
the smallest positive source-precision constant $c_*$, and a
common bound on source counts and prefactors.  The uniform
$o(U)$ condition permits a common start after which every
fixed source loss is at most $c_*U/4$.  Equation
\eqref{eq:separation-uniform-source-envelope} then supplies
a common integrable envelope $M e^{-u}$ for the remaining
negative cost.  Hence the unweighted Cauchy corrections have
norm at most $CM e^{-u_b}$ on the full forward domains.
For the weighted correction the same calculation, with the
spare half-cost retained, gives $C_1e^{K_H}$ with one $C_1$
independent both of $H$ and of the test.  This argument uses
the full-domain quantitative bounds for the fixed labels;
pointwise eventual bounds on unrelated paths would not suffice.

The proxy side buffers diverge uniformly and absorb these
absolute corrections.  The source-ratio estimates use only
the common finite real and angular flag constants, so their
constants and finite thresholds are common as well.  In the
growth argument choose one sufficiently large fixed backstep
using the common lower bound $f_p'\geq cf_p$.  Its left
boundary contribution is bounded by a common multiple of
\[
 \exp((1+\kappa)f_p(u_b)-\kappa f_p(u_a)),
\]
which tends to zero uniformly with this choice.  The possible
$D_p$ loss stays explicit; choosing one
$R=C_1\log(2+D_p)$ with $\kappa C_1>1$ makes
$D_p/\cosh(\kappa R)$ bounded uniformly.  The uniformly
available safe columns and entering collars supply the required
boundary estimates.  Their numerical locations may vary with
the test within the fixed backstep intervals; one column for a
continuum of tests is not required.  No bound on the time of return to a smaller
ordinary box has entered these calculations.

For each fixed test and $H$, its own preliminary real estimate
establishes half-domain boundedness.  The subsequent full-domain
maximum principle uses only the common side and left estimates
and the common $C_1e^{K_H}$ bound.  It therefore gives one
$C_0$ in \eqref{eq:separation-weighted} for the whole family,
regardless of the preliminary constants.  The source-ratio
lemmas then choose the same small multipliers in
$H=c_1e^{g_0}/g'_0$ or $H=e^{c_2g_0}$, and the same wide-strip
constants $W,M_0,B_1,c_4$, throughout each of the stated slope
regimes.  This proves uniformity of the upgrades.  When $n=1$,
the limit $H\to\infty$ is taken separately at each target;
a common preliminary start for all $H$ is unnecessary.

For terminal comparison, the common finite-window modulus
makes each lateral error smaller than any fixed $\eta>0$
after a common far anchor.  The Cauchy errors are uniformly
$O(e^{-u_b})$, and the fixed-window $X$ distances to both
ends tend uniformly to infinity.  The end cosine barriers
therefore disappear uniformly.  At the real targets the
affine interpolation has weights tending uniformly to $1/2$,
and bounds the difference from either terminal coefficient by
a fixed bounded constant times a positive power of $\eta$.
Letting $\eta\downarrow0$ identifies the two coefficients
uniformly there.  Once this equality is known, the errors on
both lateral boundaries are at most $\eta$ relative to their
common coefficient.  The constant lateral bound, together
with the vanishing end barriers, gives convergence throughout
the inner wedge, even for targets approaching a side where
one of the two interpolation weights tends to zero.
This proof needs slow variation near the target and
common fixed-label estimates on the full domain, but no common
time at which every deeper ordinary box is entered.
\end{proof}
